% ===== PRÉAMBULE COMMUN — À COPIER TEL QUEL EN TÊTE DE CHAQUE FICHIER .tex =====
% Ne modifier QUE :
%   - les deux lignes \fancyhead[L] et \fancyhead[R]
%   - le bloc titre \begin{center}...\end{center} juste après \begin{document}
\documentclass[11pt,a4paper]{article}
\usepackage[utf8]{inputenc}
\usepackage[T1]{fontenc}
\usepackage{lmodern}
\usepackage[french]{babel}
\usepackage{amsmath,amssymb,mathtools}
\usepackage{icomma}
\usepackage{xcolor}
\usepackage{tikz}
\usepackage{pgfplots}
\usepackage[most]{tcolorbox}
\usepackage{enumitem}
\usepackage{array,booktabs,colortbl}
\usepackage[a4paper,margin=2.2cm,headheight=26pt,headsep=10pt]{geometry}
\usepackage{fancyhdr}
\usepackage[hidelinks]{hyperref}
\usetikzlibrary{arrows.meta,calc,angles,quotes,positioning,decorations.markings,intersections,backgrounds}
\pgfplotsset{compat=1.18}

\definecolor{primary}{RGB}{226,74,88}
\definecolor{primarydark}{RGB}{150,28,42}
\definecolor{methcol}{RGB}{40,90,150}
\definecolor{excol}{RGB}{0,135,90}
\definecolor{defcol}{RGB}{0,114,178}
\definecolor{attncol}{RGB}{200,40,55}
\definecolor{thmcol}{RGB}{0,140,120}
\definecolor{courbe}{RGB}{32,110,200}
\definecolor{courbeder}{RGB}{210,70,40}
\definecolor{tang}{RGB}{0,150,90}
\definecolor{grille}{RGB}{200,205,215}
\definecolor{plus}{RGB}{200,60,60}
\definecolor{moins}{RGB}{30,90,200}

\tcbset{boxbase/.style={enhanced,breakable,boxrule=0.9pt,arc=2.5pt,
  left=3mm,right=3mm,top=2.3mm,bottom=2.3mm,fonttitle=\bfseries,coltitle=white}}
\newtcolorbox{cle}[1][]{boxbase,colback=defcol!6,colframe=defcol,
  title={\textsc{À retenir}\ifx\relax#1\relax\relax\else\textmd{ : \itshape #1}\fi}}
\newtcolorbox{methode}[1][]{boxbase,colback=methcol!7,colframe=methcol,sharp corners,
  title={\textsc{Méthode}\ifx\relax#1\relax\relax\else\textmd{ --- \itshape #1}\fi}}
\newtcolorbox{exemplebox}[1][]{boxbase,colback=excol!5,colframe=excol,
  title={\textsc{Exemple}\ifx\relax#1\relax\relax\else\textmd{ : \itshape #1}\fi}}
\newtcolorbox{piege}[1][]{boxbase,colback=attncol!6,colframe=attncol,
  title={\textsc{Attention}\ifx\relax#1\relax\relax\else\textmd{ : \itshape #1}\fi}}
\newtcolorbox{exo}[1][]{boxbase,colback=primary!4,colframe=primary,
  title={\textsc{Exercice}\ifx\relax#1\relax\relax\else\textmd{~#1}\fi}}
\newtcolorbox{corr}[1][]{boxbase,colback=thmcol!5,colframe=thmcol,
  title={\textsc{Corrigé}\ifx\relax#1\relax\relax\else\textmd{~#1}\fi}}

\newcommand{\R}{\mathbb{R}}
\newcommand{\N}{\mathbb{N}}
\newcommand{\ds}{\displaystyle}
\newcommand{\tg}{\operatorname{tg}}
\newcommand{\cotg}{\operatorname{cotg}}
\newcommand{\arctg}{\operatorname{arctg}}
\newcommand{\arccotg}{\operatorname{arccotg}}
\newcommand{\limh}{\lim_{h\to 0}}

\pagestyle{fancy}\fancyhf{}
\fancyhead[L]{\small\textcolor{primarydark}{\textbf{Thème 4} — Variations, extremums \& concavité}}
\fancyhead[R]{\small\textcolor{primarydark}{\textsc{Tutoriel}}}
\fancyfoot[C]{\small\thepage}
\renewcommand{\headrulewidth}{0.8pt}

\begin{document}

\begin{center}
{\LARGE\bfseries\color{primarydark} Variations, extremums \& concavité}\\[2pt]
{\color{primary}\rule{0.5\textwidth}{1.2pt}}\\[4pt]
{\small\itshape Lire le signe de $f'$ et de $f''$ pour comprendre la courbe}
\end{center}
\vspace{2mm}

\begin{cle}[le signe des dérivées raconte toute l'histoire de la courbe]
\textbf{Signe de $f'$ $\Longleftrightarrow$ variations.}
\begin{itemize}[leftmargin=6mm,itemsep=1pt]
  \item $f'(x)>0$ sur un intervalle $\Rightarrow$ $f$ \textbf{croissante} ($\nearrow$).
  \item $f'(x)<0$ sur un intervalle $\Rightarrow$ $f$ \textbf{décroissante} ($\searrow$).
  \item $f'$ s'annule \emph{en changeant de signe} $\Rightarrow$ \textbf{extremum} :
        $+$ puis $-$ $=$ \textbf{maximum} ; $-$ puis $+$ $=$ \textbf{minimum}.
\end{itemize}
\smallskip
\textbf{Signe de $f''$ $\Longleftrightarrow$ concavité.}
\begin{itemize}[leftmargin=6mm,itemsep=1pt]
  \item $f''(x)>0$ $\Rightarrow$ courbe \textbf{convexe} (concave vers le haut, « bol » $\smile$).
  \item $f''(x)<0$ $\Rightarrow$ courbe \textbf{concave} (concave vers le bas, « cloche » $\frown$).
  \item $f''$ s'annule \emph{en changeant de signe} $\Rightarrow$ \textbf{point d'inflexion}.
\end{itemize}
\end{cle}

\begin{methode}[dresser un tableau de variation en 6 étapes]
\begin{enumerate}[leftmargin=7mm,itemsep=1pt]
  \item \textbf{Domaine} de $f$ (valeurs interdites : dénominateur nul, $\ln$ négatif\dots).
  \item Calculer la dérivée $f'(x)$.
  \item Résoudre $f'(x)=0$ pour trouver les valeurs critiques.
  \item Étudier le \textbf{signe de $f'$} en \textbf{factorisant} (produit de facteurs, règle des signes).
  \item Reporter le signe dans le tableau et tracer les \textbf{flèches} ($\nearrow$ si $+$, $\searrow$ si $-$).
  \item Repérer les \textbf{extremums} là où $f'$ change de signe et calculer leur valeur $f(x_0)$.
\end{enumerate}
\end{methode}

\begin{piege}[une dérivée qui s'annule ne donne pas toujours un extremum]
Il faut un \textbf{changement de signe} de $f'$. Par exemple $f(x)=(3x-1)^2$ ou $f(x)=x^3$ :
la dérivée s'annule mais reste du \emph{même} signe de part et d'autre. Alors \textbf{aucun extremum}
(pour $x^3$, c'est un point d'inflexion à tangente horizontale).
\end{piege}

\begin{center}
\begin{tikzpicture}[scale=1]
% ---- Tableau de variation type : f croissante puis décroissante (max en a) ----
% cadre
\draw[thick] (0,0) rectangle (8,2.4);
\draw[thick] (0,1.6) -- (8,1.6); % ligne sous x
\draw[thick] (0,0.8) -- (8,0.8); % ligne sous f'
\draw[thick] (1.4,0) -- (1.4,2.4); % séparateur libellés
% libellés de lignes
\node at (0.7,2.0) {$x$};
\node at (0.7,1.2) {$f'(x)$};
\node at (0.7,0.4) {$f(x)$};
% valeurs de x
\node at (1.9,2.0) {$-\infty$};
\node at (4.5,2.0) {$a$};
\node at (7.6,2.0) {$+\infty$};
% signe de f'
\node[plus] at (3.0,1.2) {$+$};
\node at (4.5,1.2) {$0$};
\node[moins] at (6.2,1.2) {$-$};
% flèches de variation
\draw[->,thick,tang] (1.9,0.15) -- (4.3,0.65);
\draw[->,thick,tang] (4.7,0.65) -- (7.3,0.15);
\node at (4.5,0.9) {\footnotesize max};
\end{tikzpicture}
\hspace{6mm}
\begin{tikzpicture}[scale=0.8]
% ---- Les deux concavités ----
\draw[thick,courbe] (0,1.2) .. controls (0.6,0) and (1.4,0) .. (2,1.2);
\node at (1,-0.55) {\footnotesize $f''>0$ (bol $\smile$)};
\node at (1,1.5) {\footnotesize convexe};
\begin{scope}[xshift=3.4cm]
\draw[thick,courbeder] (0,0) .. controls (0.6,1.2) and (1.4,1.2) .. (2,0);
\node at (1,-0.55) {\footnotesize $f''<0$ (cloche $\frown$)};
\node at (1,1.5) {\footnotesize concave};
\end{scope}
\end{tikzpicture}
\end{center}

\begin{exemplebox}[minimum d'une parabole]
Soit $f(x)=x^2-4x+3$ sur $\R$.
\begin{itemize}[leftmargin=6mm,itemsep=1pt]
  \item Dérivée : $f'(x)=2x-4$.
  \item $f'(x)=0 \iff x=2$.
  \item Signe : $f'(x)<0$ pour $x<2$ ($\searrow$), $f'(x)>0$ pour $x>2$ ($\nearrow$).
  \item Donc $f'$ passe de $-$ à $+$ : \textbf{minimum} en $x=2$, de valeur $f(2)=4-8+3=-1$.
\end{itemize}
\end{exemplebox}

\end{document}
